\begin{textsample}{POS Dim 9 – global\_north\_2023\_10 – Score 1.00 – global\_north\_2023\_10/102621976.txt}  \label{ex:f9_pos_302}
The Palestinian people make up one of the largest refugee groups in the world, often forced off their land or pushed to flee amid ongoing conflict in the region, according to the United Nations.
They are stateless, their identity not defined by specific borders, but rather tied together by a collective longing for a place to call home.
Their existence has long been traded and transferred by different ruling powers throughout their history, leaving their fate hanging in the balance.
Palestinians are a diverse, multi-religious and multicultural group of roughly 14 million people internationally who trace their roots back to what is now known as the Israeli-Palestinian region of the Middle East.
Many of them no longer remain there, spread across the globe.
Those who remain in the Palestinian territory of Gaza have been under blockade, with Israel controlling nearly everything that comes in and out of the region, including food, fuel and people, and Egypt controlling Gaza ' s southern border.
Palestinians have been described as a"young society,"with people between 15 and @ @ @ @ @ @ @ @ @ @ to the United Nations Population Fund.
Ussama Makdisi, a history professor at the University of California Berkeley, told ABC News that misconceptions about Palestinian serve to dehumanize the population experiencing oppression in the Middle East, Oct. 12, 2023.
ABC News Amid the ongoing Israel-Hamas conflict in the Middle East, Palestinians say they have been"dehumanized"and stereotyped amid the international backlash to Hamas ' terrorist actions in Israel.
In Gaza, at least 1,799 people have been killed in retaliatory strikes by Israel, with an estimated 7,388 more injured, numbers that are expected to climb.
At least 1,300 people have been killed and 3,227 others have been injured in Israel from Hamas attacks.
Palestinian historians add that the people ' s long history of being ignored, attacked and displaced has left many misconceptions about the population unaddressed -- something they now hope to remedy."As soon as you identify with them as people -- as ordinary human beings -- you can not do to them what ' s now @ @ @ @ @ @ @ @ @ @ Makdisi, a history professor at the University of California Berkeley.
For example, support for the"Free Palestine"or liberation movement has historically been painted as antisemitic.
Palestinian historians say this is another misperception.
The fight for liberation, Makdisi said, is about being able to"live with dignity and equality, in a political and social and economic environment where they can thrive like any other people."According to the United Nations, 81 \% of the population in Gaza lives in poverty, with food insecurity plaguing 63 \% of Gaza citizens.
The unemployment rate is 46.6 \%, and access to clean water and electricity remains at"crisis"levels of inaccessibility, the U.N. states."Liberation means living with equality, which everybody, I think, once they understand that, most people I would say agree with it,"says Makdisi."It ' s exactly what we wish for any people in the world, any people."' Struggle for liberation ' @ @ @ @ @ @ @ @ @ @ in Palestinian territories, they love life, added Yousef Munayyer, a Palestinian-American writer and political analyst."People are looking for ways in the most difficult of moments to celebrate the beauty of life, even under oppression,"Munayyer said."This does not make Palestinians unique.
It just makes them as human as everybody else."Palestinian musicians perform in Jerusalem ' s Old City to celebrate breaking the fast on the \textbf{eighth} day of the Islamic holy month of Ramadan on May 24, 2018.
Ahmad Gharabli/AFP via Getty Images Palestinians pick olives on agricultural land in the Gaza Strip on Oct. 9, 2022.
Momen Faiz/AP All of the Palestinian scholars with whom ABC News spoke said the Palestinian passion for fighting to be able to live with dignity has stretched beyond the movement for their own people.
Palestinians have been a prominent voice in grassroots activism internationally, including the 2020 anti-racism movement that began in the United States and spread overseas.
The young population of Palestinians particularly @ @ @ @ @ @ @ @ @ @ global social movements, historians say.
Nassar believes it ' s because younger generations believe their international causes are connected, united in dismantling longstanding systems of oppression.
Palestinian resiliency The Arabic idea of"sumud,"or steadfastness, is a key part of how many Palestinians understand their collective identity, according to Maha Nassar, a history and Islamic studies professor at the University of Arizona.
It highlights the near-universal call for liberation, freedom and equality among Palestinians existing in a region under constant conflict -- a yearning for peace on Palestinian lands.
Maha Nassar, a history and Islamic studies professor at the University of Arizona, speaks with ABC News on Oct. 12, 2023.
ABC News Sumud can manifest in many ways.
It could be as simple as proclaiming one ' s pride in their Palestinian heritage, or as complicated as grassroots organizing in social movements."Palestinians around the world have taken immense pride and gone to great lengths to ensure that their Palestinian culture and identity is not @ @ @ @ @ @ @ @ @ @,"said Nassar.

% matched lemmas: eighth
\end{textsample}
